\chapter{Returning home}

Not in the literal sense of it, but directed toward all
the things you have to do after your travels. From
updating your processes based on experience to wiping
exposed devices before rejoining them to your networks.

\section{Secure your accounts and passwords}

Once you're back, if you suspect that any account
credentials you used while abroad might have been exposed
(especially if you had to log in over a hotel or
conference Wi-Fi), address the issue. Change passwords for
any accounts you accessed unsafely during the trip -- but
if it doesn't seem necessary, be mindful that changing
them unnecessarily can sometimes introduce additional
risk.

Do this from a trusted device/network (ideally wait until
you're on your home or office network, not the airport
Wi-Fi). Use this opportunity to upgrade weak passwords and
ensure 2FA is still working on those accounts.

Check your email filters and crypto account settings for
any unauthorized changes (attackers sometimes add
forwarding rules or new withdrawal addresses if they did
get access).

Essentially, rotate secrets that may have been used under
less-secure conditions.

\section{Inspect and clean your devices}

After traveling, give your devices a thorough once-over.
Run a reputable anti-malware scan on laptops and phones.
Look for any unusual apps, processes, or device behavior
(for example, unusual battery drain could indicate
malware).

If you were in a high-risk environment or your device was
out of your control at any point, consider wiping the
device and restoring it from your pre-trip backup (or
factory-resetting a phone) to ensure it's
clean.\footnote{\url{https://architectsecurity.org/2017/10/mobile-device-security-for-international-travelers-part-3-how-to-clean-up-your-mobile-devices-after-international-travel}}
This is especially recommended for ``burner'' devices used
on the trip -- you can safely restore your data onto your
main device and decommission the travel device.

For hardware wallets, verify they weren't tampered with:
check the device seals if any, and when you connect,
confirm the firmware is still legitimate (if the
manufacturer provides verification software). If you have
any suspicion that a device (or hardware wallet) was
compromised, do not continue using it for sensitive
transactions. Transfer crypto assets to new wallets (using
your seed backups in a new device if necessary) once
you're on a secure network and device.

It's also a good idea to disable or remove any
travel-specific eSIMs or accounts you used on the trip --
for example, remove that foreign cellular plan from your
phone if you no longer need it, and uninstall any travel
or conference apps that are no longer required.

\section{Post-travel review}

Now that you're home, reflect on the trip and note any
security incidents or close calls. If any device was lost,
stolen, or even taken out of your sight and potentially
tampered (like held by airport security for a long
inspection), inform your organization's IT or security
team immediately. They may assist with forensic checks or
account monitoring. Also inform colleagues if any work
data might have been exposed so they can be vigilant. This
is not about blame -- it's about mitigating any damage
early.

Re-enable any data or accounts you put in ``travel mode.''
For instance, if you used 1Password Travel Mode to hide
vaults, log in and turn those vaults back on. If you
created throwaway emails or burner chat accounts for the
trip, decide if you'll deactivate them now.

\subsection{Update your threat model}

Based on your experience: did any new threats emerge, or
did some precautions prove unnecessary? Use that to
improve future travel prep. Finally, share key lessons
with your team. Sharing what you've learned from each trip
and tweaking your security practices contributes to a
stronger security culture for everyone!

There's more information for high profile individuals in
the following chapter.
